\documentclass[11pt]{article}

\usepackage[margin=1in]{geometry}
\usepackage{amsmath,amssymb,amsthm}
\usepackage{graphicx}
\usepackage{hyperref}
\usepackage{enumitem}

\hypersetup{
  colorlinks=true,
  linkcolor=blue,
  urlcolor=blue,
  citecolor=blue
}

\newtheorem{theorem}{Theorem}
\newtheorem{lemma}{Lemma}
\theoremstyle{remark}
\newtheorem{remark}{Remark}

\newcommand{\D}{\mathbb D}
\newcommand{\T}{\mathbb T}
\newcommand{\N}{\mathbb N}
\newcommand{\C}{\mathbb C}

\title{Spaceability of Smooth Disk-Algebra Functions with Boundary Pringsheim Singularities}
\author{Math discovery smoke-test packet}
\date{June 12, 2026}

\begin{document}
\maketitle

\section{Source}

The source paper is

\begin{quote}
Andreas Debrouwere, \emph{Some results about spaceability in function spaces},
arXiv:2606.13195.
\end{quote}

Remark 3.15 points to the following question from Bernal-Gonzalez et al.:

\begin{center}
\includegraphics[width=\textwidth]{figures/open_problem_crop.png}
\end{center}

Here \(A^\infty(\D)\) denotes the smooth disk algebra.  A boundary point is
Pringsheim singular for the boundary restriction when the holomorphic
function has no analytic continuation through that point.  Thus it is enough
to construct a closed infinite-dimensional subspace of \(A^\infty(\D)\) whose
nonzero elements all have \(\T\) as a natural boundary.

\section{Claim}

\begin{theorem}
The set
\[
\{f\in A^\infty(\D): f|_{\T}\text{ is Pringsheim singular at every boundary point}\}
\]
is spaceable in \(A^\infty(\D)\).
\end{theorem}

\section{Proof}

We use the classical Hadamard gap theorem in the following form.

\begin{lemma}[Hadamard gap theorem]
Let
\[
F(z)=\sum_{\ell=1}^\infty a_\ell z^{m_\ell}
\]
have radius of convergence \(1\), where \(m_{\ell+1}/m_\ell\ge q>1\).
Then the unit circle is a natural boundary for \(F\).
\end{lemma}

Let the master lacunary sequence be \(N_k=2^k\), \(k\ge 1\).  Partition
\(\N\) into infinitely many infinite sets
\[
A_m=\{2^m(2j+1): j=0,1,2,\ldots\},\qquad m=0,1,2,\ldots .
\]
For \(m\ge0\), define
\[
n_{m,j}=2^{2^m(2j+1)}
\]
and
\[
g_m(z)=\sum_{j=0}^{\infty}\exp(-\sqrt{n_{m,j}})\,z^{n_{m,j}}.
\]
The sets \(\{n_{m,j}:j\ge0\}\) are pairwise disjoint and all lie in the
master support \(\{2^k:k\ge1\}\).  Moreover,
\[
\sum_j n_{m,j}^r \exp(-\sqrt{n_{m,j}})<\infty
\]
for every \(r\ge0\).  Hence each derivative series of \(g_m\) converges
uniformly on \(\overline{\D}\), so \(g_m\in A^\infty(\D)\).

Let
\[
X=\overline{\operatorname{span}\{g_m:m=0,1,2,\ldots\}}^{A^\infty(\D)}.
\]
This is a closed infinite-dimensional subspace, since the \(g_m\)'s have
pairwise disjoint nonzero Taylor supports and are therefore linearly
independent.

We claim that every nonzero \(f\in X\) has \(\T\) as a natural boundary.
Write
\[
f(z)=\sum_{n=0}^{\infty} c_n z^n .
\]
The coefficient functionals \(f\mapsto c_n=f^{(n)}(0)/n!\) are continuous on
\(A^\infty(\D)\).  Since finite linear combinations of the \(g_m\)'s have no
coefficients outside the master set, every \(f\in X\) also satisfies
\[
c_n=0\qquad(n\notin\{2^k:k\ge1\}).
\]
Likewise, for fixed \(m\), finite linear combinations satisfy
\[
\frac{c_{n_{m,j}}}{\exp(-\sqrt{n_{m,j}})}
=
\frac{c_{n_{m,j'}}}{\exp(-\sqrt{n_{m,j'}})}
\qquad(j,j'\ge0),
\]
and this equality is preserved under closure.  Therefore each \(f\in X\)
has scalars \(\lambda_m\) such that
\[
c_{n_{m,j}}=\lambda_m\exp(-\sqrt{n_{m,j}})\qquad(j\ge0).
\]
If \(f\ne0\), then \(\lambda_m\ne0\) for at least one \(m\); otherwise all
Taylor coefficients of \(f\) would be zero.

For such an \(m\),
\[
\lim_{j\to\infty}
|c_{n_{m,j}}|^{1/n_{m,j}}
=
\lim_{j\to\infty}
|\lambda_m|^{1/n_{m,j}}\exp(-1/\sqrt{n_{m,j}})
=1.
\]
Since \(f\in A^\infty(\D)\), its radius of convergence is at least \(1\).
The preceding limit shows that the radius is exactly \(1\).  The nonzero
Taylor support of \(f\) is a subsequence of the master support \(\{2^k\}\);
therefore, when listed increasingly, its successive exponents have ratio at
least \(2\).  Hadamard's gap theorem applies and gives \(\T\) as a natural
boundary for \(f\).  Thus \(f|_{\T}\) is Pringsheim singular at every
boundary point.

Therefore \(X\setminus\{0\}\) is contained in the set in the theorem, and
the set is spaceable.
\qed

\section{Verification Notes}

The script \texttt{code/check\_lacunary\_construction.py} checks finite
samples of the partition and coefficient profile:

\begin{verbatim}
conda run --no-capture-output -n sandbox python code/check_lacunary_construction.py
\end{verbatim}

The run confirmed disjoint blocks, a master consecutive ratio at least \(2\),
and the expected root behavior of \(\exp(-\sqrt n)\).  These checks are only
sanity checks; the proof above is analytic and uses Hadamard's theorem.

\section{Limitations and Review Focus}

The main external ingredient is the classical Hadamard gap theorem.  A human
reviewer should verify that the terminology ``Pringsheim singular at every
point'' in Bernal-Gonzalez et al. agrees with the natural-boundary conclusion
used here.  The construction gives the stronger conclusion that no nonzero
element of \(X\) analytically continues through any boundary point.

\end{document}
